\subsection{Matter Perturbations and Structure Growth}\label{subsec:ps_matter}

This subsection describes the matter power spectrum and how it evolves with cosmic time.

\subsubsection{Linear perturbation theory}

At early times, density perturbations $\delta(\mathbf{x}, z) = (\rho - \bar{\rho})/\bar{\rho}$ are small and grow linearly:

\begin{equation}
\delta(k, z) = D_+(z) \delta(k, z=0)
\end{equation}

where $D_+(z)$ is the linear growth factor, normalized to 1 today.

The matter power spectrum evolves as:

\begin{equation}
P(k, z) = D_+^2(z) \, P(k, z=0) = D_+^2(z) \, P_0(k)
\end{equation}

where $P_0(k)$ is a reference spectrum (usually at $z=0$ or early times).

\subsubsection{Linear vs nonlinear scales}

- **Linear scales** ($k < k_{\text{nl}}$): Perturbation theory accurate; $P(k) \propto D_+^2$
- **Nonlinear scales** ($k > k_{\text{nl}}$): Mode coupling; power enhanced; requires $N$-body simulations or fitting formulas

The nonlinear scale $k_{\text{nl}}(z)$ increases with redshift (less time for collapse):

\begin{equation}
k_{\text{nl}}(z) \approx 0.12 \, h \, \text{Mpc}^{-1} \quad (z=0) \quad \rightarrow \quad 0.25 \, h \, \text{Mpc}^{-1} \quad (z=2)
\end{equation}

\subsubsection{Galaxy bias}

Galaxies do not trace the matter field 1-to-1. Linear galaxy bias relates galaxy and matter power spectra:

\begin{equation}
P_{\text{gal}}(k, z) = b^2(z) \, P_{\text{matter}}(k, z)
\end{equation}

where $b(z)$ is the linear bias parameter ($b > 1$ for rare, massive objects; $b \approx 1$ for typical galaxies).

\example{
ΛCDM with $\Om=0.3, \OL=0.7$ gives:
- At $z=0$: $D_+(0) \approx 1$ (normalized)
- At $z=1$: $D_+(1) \approx 0.76$ (structure grows more slowly at high-$z$)
- Typical galaxy bias: $b \approx 1.5$--$2$ at $z \sim 1$
}

See \cref{subsec:background_evolution} for growth factor evolution and \cref{ch:weak_lensing} for weak lensing kernels involving matter power.
